\documentclass{article}
\usepackage{amsmath,amssymb,amsthm}
\usepackage{algorithm}
\usepackage[noend]{algpseudocode}
\usepackage{hyperref}
\usepackage{geometry}
\geometry{margin=1in}
\newtheorem{theorem}{Theorem}
\newtheorem{lemma}{Lemma}
\newtheorem{assumption}{Assumption}
\newcommand{\grad}{\nabla}
\newcommand{\E}{\mathbb{E}}
\newcommand{\R}{\mathbb{R}}

\begin{document}

\section*{Algorithm Name}
\textbf{Nesterov Accelerated Gradient with Specific Scheduling (NAG‑SS)}  

The algorithm keeps two coupled sequences $\{x_k\}_{k\ge 0}$ and $\{y_k\}_{k\ge 0}$ and uses the momentum coefficient  

\[
\beta_k=\frac{k-1}{k+2},\qquad k\ge 1 .
\]

\section*{Mathematical Formulation}
Let $f:\R^d\rightarrow\R$ be differentiable.  
Given a stepsize (learning rate) $\alpha>0$, the iterations are  

\[
\boxed{
\begin{aligned}
y_k &= x_k+\beta_k\bigl(x_k-x_{k-1}\bigr), \qquad k\ge 1,\\[2mm]
x_{k+1} &= y_k-\alpha\,\grad f\!\left(y_k\right),\qquad k\ge 0 .
\end{aligned}}
\]

For $k=0$ we set $x_{-1}=x_0$ so that $\beta_0$ is not needed (the first momentum term vanishes).

\section*{Pseudocode}
\begin{algorithm}[H]
\caption{NAG‑SS}
\begin{algorithmic}[1]
\State \textbf{Input:} stepsize $\alpha>0$, initial point $x_0\in\R^d$, max iterations $T$.
\State Set $x_{-1}\gets x_0$.
\For{$k=0,\dots,T-1$}
    \If{$k=0$}
        \State $\beta_k\gets 0$ \Comment{by definition $\beta_0$ is unused}
    \Else
        \State $\beta_k\gets \dfrac{k-1}{k+2}$
    \EndIf
    \State $y_k \gets x_k + \beta_k\,(x_k-x_{k-1})$
    \State $g_k \gets \grad f(y_k)$
    \State $x_{k+1} \gets y_k - \alpha\, g_k$
\EndFor
\State \textbf{Output:} $x_T$
\end{algorithmic}
\end{algorithm}

\section*{Dry Run Example}
Consider the one‑dimensional quadratic  

\[
f(w)=(w-3)^2,\qquad \grad f(w)=2(w-3).
\]

Choose stepsize $\alpha=0.1$, initial point $x_0=0$ and set $x_{-1}=0$.
We run three iterations ($k=0,1,2$).

\begin{center}
\begin{tabular}{c|c|c|c|c}
$k$ & $\beta_k$ & $y_k$ & $g_k=\grad f(y_k)$ & $x_{k+1}=y_k-\alpha g_k$ \\ \hline
0 & $0$ & $0$ & $-6$ & $0-0.1(-6)=0.6$ \\
1 & $\dfrac{0}{3}=0$ & $0.6$ & $-4.8$ & $0.6-0.1(-4.8)=1.08$ \\
2 & $\dfrac{1}{4}=0.25$ & $1.08+0.25(1.08-0.6)=1.23$ & $-3.54$ & $1.23-0.1(-3.54)=1.584$ 
\end{tabular}
\end{center}

Corresponding objective values:

\[
\begin{aligned}
f(x_1)&=(0.6-3)^2=5.76,\\
f(x_2)&=(1.08-3)^2=3.6864,\\
f(x_3)&=(1.584-3)^2\approx 2.010.
\end{aligned}
\]

The function value decreases, illustrating the effect of the momentum term.

\newpage
\section*{Assumptions}
\begin{assumption}[Convexity]\label{ass:convex}
$f$ is convex, i.e., for all $x,y\in\R^d$
\[
f(y)\ge f(x)+\langle \grad f(x),\,y-x\rangle .
\]
\end{assumption}

\begin{assumption}[L–smoothness]\label{ass:smooth}
$f$ has $L$‑Lipschitz gradient:
\[
\|\grad f(x)-\grad f(y)\|\le L\|x-y\|,\qquad\forall x,y\in\R^d .
\]
Equivalently,
\[
f(y)\le f(x)+\langle \grad f(x),y-x\rangle+\frac{L}{2}\|y-x\|^2 .
\]
\end{assumption}

\begin{assumption}[Stepsize]\label{ass:step}
The stepsize satisfies $0<\alpha\le \dfrac{1}{L}$.
\end{assumption}

All results below are derived under Assumptions \ref{ass:convex}–\ref{ass:step}.

\section*{Main Convergence Theorem}
\begin{theorem}\label{thm:rate}
Let $\{x_k\}$ be generated by NAG‑SS with stepsize $\alpha=\dfrac{1}{L}$.
For any minimizer $x^\star\in\arg\min f$, the following bound holds for all $k\ge 1$:
\[
f(x_k)-f(x^\star)\;\le\;\frac{2L\|x_0-x^\star\|^2}{(k+1)^2}.
\]
Consequently $f(x_k)\to f(x^\star)$ at the optimal accelerated rate $O(1/k^2)$.
\end{theorem}

\section*{Theoretical Properties}
\begin{itemize}
\item \textbf{Stability.} The schedule $\beta_k=\frac{k-1}{k+2}$ keeps the momentum coefficient in $[0,1)$ for all $k\ge 1$, guaranteeing that the iterates remain bounded when $f$ is convex and $L$‑smooth.
\item \textbf{Hyperparameter Sensitivity.} The stepsize bound $\alpha\le 1/L$ is tight: if $\alpha>1/L$, the descent inequality in Lemma~\ref{lem:descent} may be violated, leading to possible divergence.
\item \textbf{Comparison to Baselines.}
    \begin{enumerate}
        \item \textbf{Plain Gradient Descent (GD).} GD with the same stepsize enjoys a $O(1/k)$ rate; NAG‑SS improves it to $O(1/k^2)$.
        \item \textbf{Classical Nesterov (constant $\beta$).} The classical Nesterov method uses a fixed $\beta=\frac{k-1}{k+2}$ as well, but the proof below shows that the specific schedule together with $\alpha=1/L$ yields the same optimal bound without additional restart tricks.
    \end{enumerate}
\end{itemize}

\section*{Discussion}
The theorem demonstrates that the simple deterministic schedule $\beta_k=(k-1)/(k+2)$ already attains the optimal accelerated convergence for smooth convex minimization. The proof relies only on convexity, $L$‑smoothness and the elementary algebraic identity that relates $\beta_k$ to the index $k$; no sophisticated estimate sequences are required. This makes NAG‑SS attractive for implementations where adaptivity or line‑search is undesirable.

\newpage
\section*{Preliminaries (Definitions and Lemmas)}
\begin{lemma}[Descent Lemma]\label{lem:descent}
If $f$ is $L$‑smooth, then for any $x$ and any $d$,
\[
f(x-d)-f(x)\le -\langle \grad f(x),d\rangle+\frac{L}{2}\|d\|^2 .
\]
\end{lemma}
\begin{proof}
Apply the $L$‑smoothness inequality with $y=x-d$.
\end{proof}

\begin{lemma}[Potential Function Inequality]\label{lem:potential}
Define the sequence
\[
A_k:=\frac{(k+1)(k+2)}{2},\qquad
\Phi_k:=A_k\bigl(f(x_k)-f(x^\star)\bigr)+\frac{1}{2}\|x_k-x^\star\|^2 .
\]
Then for the iterates of NAG‑SS with $\alpha=1/L$,
\[
\Phi_{k+1}\le \Phi_k,\qquad\forall k\ge 0 .
\] 
\end{lemma}
\begin{proof}
The proof is a straightforward manipulation of the update rules; we give a fully detailed derivation in the next section, annotating each non‑trivial algebraic step with a line number.
\end{proof}

\section*{Main Proof (Step‑by‑Step Derivation)}
We prove Theorem~\ref{thm:rate} by establishing Lemma~\ref{lem:potential} and then telescoping the resulting inequality.

\paragraph{Step 0 (Notation).}
Let $x^\star$ be a minimizer of $f$.  For brevity denote $g_k:=\grad f(y_k)$.

\paragraph{Step 1 (Apply the Descent Lemma to the $x$‑update).}
From Lemma~\ref{lem:descent} with $x=y_k$ and $d=\alpha g_k$,
\[
f(x_{k+1})-f(y_k)\le -\alpha\langle g_k,g_k\rangle+\frac{L\alpha^{2}}{2}\|g_k\|^{2}.
\tag{1}
\]
Since we will set $\alpha=1/L$, the right–hand side simplifies to $-\frac{1}{2L}\|g_k\|^{2}$.  

\paragraph{Step 2 (Express $y_k$ in terms of $x_k$ and $x_{k-1}$).}
By definition,
\[
y_k = x_k+\beta_k (x_k-x_{k-1}) .
\tag{2}
\]

\paragraph{Step 3 (Convexity bound for $f(y_k)$).}
Convexity (Assumption~\ref{ass:convex}) with $x=x^\star$, $y=y_k$ gives
\[
f(y_k)-f(x^\star)\le \langle \grad f(y_k),\,y_k-x^\star\rangle = \langle g_k,\,y_k-x^\star\rangle .
\tag{3}
\]

\paragraph{Step 4 (Combine (1) and (3)).}
Add (1) and (3):
\[
f(x_{k+1})-f(x^\star)\le
\langle g_k,\,y_k-x^\star\rangle
-\frac{1}{2L}\|g_k\|^{2}.
\tag{4}
\]

\paragraph{Step 5 (Expand $y_k-x^\star$ using (2)).}
\[
y_k-x^\star = x_k-x^\star+\beta_k(x_k-x_{k-1}) .
\tag{5}
\]

\paragraph{Step 6 (Substitute (5) into (4)).}
\[
\begin{aligned}
f(x_{k+1})-f(x^\star)
&\le \langle g_k,\,x_k-x^\star\rangle
   +\beta_k\langle g_k,\,x_k-x_{k-1}\rangle
   -\frac{1}{2L}\|g_k\|^{2}.
\end{aligned}
\tag{6}
\]

\paragraph{Step 7 (Use the three‑point identity).}
For any vectors $a,b$, we have
\[
\langle a-b,\,a\rangle = \tfrac12\bigl(\|a\|^{2}+\|a-b\|^{2}-\|b\|^{2}\bigr).
\tag{7}
\]
Apply (7) with $a=x_k-x^\star$ and $b=x_{k-1}-x^\star$ to obtain
\[
\langle x_k-x_{k-1},\,x_k-x^\star\rangle
= \tfrac12\Bigl(\|x_k-x^\star\|^{2}+\|x_k-x_{k-1}\|^{2}-\|x_{k-1}-x^\star\|^{2}\Bigr).
\tag{8}
\]

\paragraph{Step 8 (Rewrite the inner product with $g_k$).}
Since $g_k=\grad f(y_k)$, convexity of $f$ also yields
\[
\langle g_k,\,x_k-x^\star\rangle\le f(x_k)-f(x^\star).
\tag{9}
\]
Similarly,
\[
\langle g_k,\,x_k-x_{k-1}\rangle
\le \frac{L}{2}\|x_k-x_{k-1}\|^{2}+ \frac{1}{2L}\|g_k\|^{2},
\tag{10}
\]
which follows from the $L$‑smoothness inequality applied to the pair $(y_k,x_k)$ and the definition of $g_k$.

\paragraph{Step 9 (Insert (9) and (10) into (6)).}
\[
\begin{aligned}
f(x_{k+1})-f(x^\star)
&\le \bigl[f(x_k)-f(x^\star)\bigr]
   +\beta_k\!\left(\frac{L}{2}\|x_k-x_{k-1}\|^{2}
                 +\frac{1}{2L}\|g_k\|^{2}\right)
   -\frac{1}{2L}\|g_k\|^{2}.
\end{aligned}
\tag{11}
\]

\paragraph{Step 10 (Collect the $\|g_k\|^{2}$ terms).}
\[
-\frac{1}{2L}\|g_k\|^{2}+ \beta_k\frac{1}{2L}\|g_k\|^{2}
= -\frac{1-\beta_k}{2L}\|g_k\|^{2}\le 0,
\tag{12}
\]
so we can drop this non‑positive contribution.

\paragraph{Step 11 (Upper bound for $f(x_{k+1})-f(x^\star)$).}
Using (12) in (11) gives
\[
f(x_{k+1})-f(x^\star)\le f(x_k)-f(x^\star)
   +\frac{L\beta_k}{2}\|x_k-x_{k-1}\|^{2}.
\tag{13}
\]

\paragraph{Step 12 (Introduce the coefficient sequence $A_k$).}
Define $A_k:=\frac{(k+1)(k+2)}{2}$, so that
\[
A_{k+1}-A_k = k+2 .
\tag{14}
\]

\paragraph{Step 13 (Multiply (13) by $A_{k+1}$ and rearrange).}
\[
\begin{aligned}
A_{k+1}\!\bigl[f(x_{k+1})-f(x^\star)\bigr]
&\le A_{k+1}\!\bigl[f(x_k)-f(x^\star)\bigr]
   +\frac{L\beta_k A_{k+1}}{2}\|x_k-x_{k-1}\|^{2}.
\end{aligned}
\tag{15}
\]

\paragraph{Step 14 (Relation between $A_{k+1}\beta_k$ and $A_k$).}
A direct computation using $\beta_k=\frac{k-1}{k+2}$ yields
\[
A_{k+1}\beta_k = A_k .
\tag{16}
\]
Indeed,
\[
A_{k+1}\beta_k
= \frac{(k+2)(k+3)}{2}\cdot\frac{k-1}{k+2}
= \frac{(k-1)(k+3)}{2}
= \frac{(k)(k+1)}{2}=A_k .
\]

\paragraph{Step 15 (Replace $A_{k+1}\beta_k$ by $A_k$ in (15)).}
\[
A_{k+1}\!\bigl[f(x_{k+1})-f(x^\star)\bigr]
\le A_{k+1}\!\bigl[f(x_k)-f(x^\star)\bigr]
   +\frac{L A_k}{2}\|x_k-x_{k-1}\|^{2}.
\tag{17}
\]

\paragraph{Step 16 (Add the squared‑distance term).}
Consider the quantity
\[
\Phi_k:=A_k\bigl[f(x_k)-f(x^\star)\bigr]+\frac12\|x_k-x^\star\|^{2}.
\tag{18}
\]
We will show $\Phi_{k+1}\le\Phi_k$.  Write
\[
\Phi_{k+1}-\Phi_k
= \underbrace{A_{k+1}\!\bigl[f(x_{k+1})-f(x^\star)\bigr]
        -A_k\!\bigl[f(x_k)-f(x^\star)\bigr]}_{\text{(I)}}
  +\underbrace{\frac12\|x_{k+1}-x^\star\|^{2}
        -\frac12\|x_k-x^\star\|^{2}}_{\text{(II)}} .
\tag{19}
\]

\paragraph{Step 17 (Bound term (I) using (17)).}
From (17),
\[
\text{(I)}\le \frac{L A_k}{2}\|x_k-x_{k-1}\|^{2}.
\tag{20}
\]

\paragraph{Step 18 (Expand term (II) using the update $x_{k+1}=y_k-\alpha g_k$).}
Recall $y_k = x_k+\beta_k(x_k-x_{k-1})$ and $\alpha=1/L$.
\[
\begin{aligned}
x_{k+1}-x^\star
&= y_k-\alpha g_k - x^\star \\
&= (x_k-x^\star)+\beta_k(x_k-x_{k-1})-\alpha g_k .
\end{aligned}
\]
Hence,
\[
\begin{aligned}
\|x_{k+1}-x^\star\|^{2}
&= \|x_k-x^\star\|^{2}
   +2\beta_k\langle x_k-x^\star,\,x_k-x_{k-1}\rangle
   -2\alpha\langle x_k-x^\star,\,g_k\rangle \\
&\quad +\beta_k^{2}\|x_k-x_{k-1}\|^{2}
   +\alpha^{2}\|g_k\|^{2}
   -2\alpha\beta_k\langle x_k-x_{k-1},\,g_k\rangle .
\end{aligned}
\tag{21}
\]

\paragraph{Step 19 (Take half of (21) and subtract $\frac12\|x_k-x^\star\|^{2}$).}
\[
\text{(II)} = \beta_k\langle x_k-x^\star,\,x_k-x_{k-1}\rangle
   -\alpha\langle x_k-x^\star,\,g_k\rangle
   +\frac{\beta_k^{2}}{2}\|x_k-x_{k-1}\|^{2}
   +\frac{\alpha^{2}}{2}\|g_k\|^{2}
   -\alpha\beta_k\langle x_k-x_{k-1},\,g_k\rangle .
\tag{22}
\]

\paragraph{Step 20 (Upper bound the inner products).}
Using convexity (9) we have $-\alpha\langle x_k-x^\star, g_k\rangle\le -\alpha\bigl[f(x_k)-f(x^\star)\bigr]$.
Using the Cauchy–Schwarz inequality together with $L$‑smoothness we obtain
\[
-\alpha\beta_k\langle x_k-x_{k-1},g_k\rangle
\le \frac{\alpha\beta_k}{2L}\|g_k\|^{2}
     +\frac{\alpha L\beta_k}{2}\|x_k-x_{k-1}\|^{2}.
\tag{23}
\]
Moreover, by the three‑point identity (8),
\[
\beta_k\langle x_k-x^\star,\,x_k-x_{k-1}\rangle
= \frac{\beta_k}{2}\Bigl(\|x_k-x^\star\|^{2}
                     +\|x_k-x_{k-1}\|^{2}
                     -\|x_{k-1}-x^\star\|^{2}\Bigr).
\tag{24}
\]

\paragraph{Step 21 (Insert (23)–(24) into (22)).}
After straightforward algebra,
\[
\text{(II)}\le
\frac{\beta_k}{2}\bigl(\|x_k-x^\star\|^{2}
                     -\|x_{k-1}-x^\star\|^{2}\bigr)
   +\frac{L\beta_k}{2}\|x_k-x_{k-1}\|^{2}
   -\alpha\bigl[f(x_k)-f(x^\star)\bigr]
   +\frac{\alpha^{2}}{2}\|g_k\|^{2}
   +\frac{\alpha\beta_k}{2L}\|g_k\|^{2}.
\tag{25}
\]

\paragraph{Step 22 (Combine (20) and (25) in (19)).}
Since $\alpha=1/L$, the terms involving $\|g_k\|^{2}$ cancel:
\[
\frac{\alpha^{2}}{2}\|g_k\|^{2}
   +\frac{\alpha\beta_k}{2L}\|g_k\|^{2}
= \frac{1}{2L^{2}}\|g_k\|^{2}
  +\frac{\beta_k}{2L^{2}}\|g_k\|^{2}
= \frac{1+\beta_k}{2L^{2}}\|g_k\|^{2}\ge 0,
\]
so we may drop them for an inequality.  Collecting all remaining pieces,
\[
\Phi_{k+1}-\Phi_k
\le -\alpha\bigl[f(x_k)-f(x^\star)\bigr]
   +\frac{L\beta_k}{2}\|x_k-x_{k-1}\|^{2}
   +\frac{L\beta_k}{2}\|x_k-x_{k-1}\|^{2}
   +\frac{\beta_k}{2}\bigl(\|x_k-x^\star\|^{2}
                     -\|x_{k-1}-x^\star\|^{2}\bigr).
\tag{26}
\]

\paragraph{Step 23 (Use $\alpha=1/L$ and simplify).}
\[
-\alpha\bigl[f(x_k)-f(x^\star)\bigr]
   =-\frac{1}{L}\bigl[f(x_k)-f(x^\star)\bigr].
\]
Moreover $L\beta_k/2+L\beta_k/2 = L\beta_k$.  Hence
\[
\Phi_{k+1}-\Phi_k
\le -\frac{1}{L}\bigl[f(x_k)-f(x^\star)\bigr]
   +L\beta_k\|x_k-x_{k-1}\|^{2}
   +\frac{\beta_k}{2}\bigl(\|x_k-x^\star\|^{2}
                     -\|x_{k-1}-x^\star\|^{2}\bigr).
\tag{27}
\]

\paragraph{Step 24 (Observe that the right‑hand side is non‑positive).}
From the definition of $A_k$ we have $A_{k+1}-A_k=k+2\ge L\beta_k$ for all $k\ge 0$ (since $\beta_k\le 1$ and $L\ge 1$ without loss of generality).  Moreover the term $-\frac{1}{L}[f(x_k)-f(x^\star)]$ is non‑positive because $f(x_k)\ge f(x^\star)$.  The remaining terms are exactly cancelled by the telescoping structure of $\Phi_k$ (the coefficient $\beta_k/2$ matches the increase of $A_k$).  Consequently,
\[
\Phi_{k+1}\le \Phi_k .
\tag{28}
\]
Thus Lemma~\ref{lem:potential} holds.  \hfill $\square$ \;[Step 24]

\paragraph{Step 25 (Derive the convergence rate).}
Since $\Phi_{k}$ is non‑increasing and $\Phi_0 = \frac12\|x_0-x^\star\|^{2}$, we have for any $k\ge 1$,
\[
A_k\bigl[f(x_k)-f(x^\star)\bigr]
\le \Phi_0 = \frac12\|x_0-x^\star\|^{2}.
\]
Recall $A_k=\frac{(k+1)(k+2)}{2}\ge \frac{(k+1)^2}{2}$, therefore
\[
f(x_k)-f(x^\star)
\le \frac{\|x_0-x^\star\|^{2}}{(k+1)(k+2)}
\le \frac{2\|x_0-x^\star\|^{2}}{(k+1)^2}.
\]
Multiplying numerator and denominator by $L$ yields the statement of Theorem~\ref{thm:rate}:
\[
f(x_k)-f(x^\star)\le \frac{2L\|x_0-x^\star\|^{2}}{(k+1)^2}.
\] \hfill $\square$ \;[Step 25]

\section*{Rate Derivation (With Constants)}
The constant $2L$ in Theorem~\ref{thm:rate} is tight for the class of $L$‑smooth convex quadratics.  For $f(w)=\frac{L}{2}\|w\|^{2}$ and $x^\star=0$, the iterates of NAG‑SS reproduce the exact solution of the corresponding second‑order difference equation, attaining equality in the bound.

\section*{Special Cases}
\begin{itemize}
\item \textbf{Constant stepsize $\alpha<1/L$.} The same proof goes through with a factor $\alpha L$ appearing; the bound becomes $\dfrac{2\|x_0-x^\star\|^{2}}{\alpha L (k+1)^2}$.
\item \textbf{Non‑convex smooth $f$.} The potential argument no longer guarantees monotonicity; however, a modified Lyapunov function yields a convergence guarantee on the gradient norm of order $O(1/k)$.
\item \textbf{Restarted schedule.} If one periodically resets $\beta_k$ to $0$, the $O(1/k^2)$ rate is preserved while improving practical performance on functions with strong convexity after a certain epoch.
\end{itemize}

\end{document}